\documentclass[a4paper,11pt]{article}
\usepackage[utf8]{inputenc}
\usepackage[T1]{fontenc}
\usepackage[french]{babel}
\usepackage[left=1cm,right=1cm,top=.5cm,bottom=0cm]{geometry}
\usepackage{enumitem}
\usepackage{hyperref}
\usepackage{xcolor}
\usepackage{titlesec}
\usepackage{fontawesome5}
\usepackage{lmodern}
\usepackage[normalem]{ulem}

\definecolor{primary}{RGB}{0, 82, 165}
\definecolor{secondary}{RGB}{80, 80, 80}
\hypersetup{colorlinks=true, linkcolor=primary, filecolor=primary, urlcolor=primary}
\titleformat{\section}{\large\bfseries\color{primary}\uppercase}{}{0em}{}[\titlerule]
\titlespacing{\section}{0pt}{8pt}{5pt}
\newcommand{\resumeItem}[1]{\item\small{#1}}

\begin{document}

\begin{center}
    {\Large \textbf{Lo\"ic Aron MBASSI EWOLO}} \\ \vspace{4pt}
    \textbf{\large Ing\'enieur MLOps / Data Scientist} \\
    \textit{\small Disponible d\`es septembre 2026 pour alternance 12 mois} \\ \vspace{4pt}
    \small
    \faMapMarker\ Cergy, France (Mobile IDF) \quad
    \faPhone\ (+33) 7 59 52 33 98 \quad
    \faEnvelope\ \href{mailto:loicaron.mbassiewolo@ensea.fr}{loicaron.mbassiewolo@ensea.fr} \\ \vspace{2pt}
    \faLinkedin\ \href{https://www.linkedin.com/in/loic-mbassi/}{linkedin.com/in/loic-mbassi} \quad
    \faGithub\ \href{https://github.com/Nameless0l}{github.com/Nameless0l} \quad
    \faGlobe\ \href{https://mbassiloic.tech}{mbassiloic.tech}
\end{center}

\vspace{-6pt}

\noindent\small{L'industrialisation de mod\`eles ML en production n\'ecessite une ma\^itrise transverse du cycle de vie des donn\'ees, de l'entra\^inement \`a la supervision en production. Au MACIF, rejoindre l'\'equipe Data \& IA pour fiabiliser des pipelines de classification documentaire et d'extraction d'information repr\'esente une opportunit\'e directement align\'ee avec mon parcours \`a l'interface entre recherche ML et d\'eploiement industriel.}

\section{Formation}
\begin{itemize}[leftmargin=*, label=\tiny$\bullet$, nosep]
    \item \textbf{ENSEA -- \'Ecole Nationale Sup\'erieure de l'\'Electronique et de ses Applications} \hfill \textit{2025 -- 2027} \\
    \small{Double dipl\^ome d'Ing\'enieur -- Syst\`emes Embarqu\'es, Signal, IA.}
    \item \textbf{\'Ecole Polytechnique de Yaound\'e (1\textsuperscript{\`ere} \'ecole d'ing\'enieur CEMAC)} \hfill \textit{2021 -- 2025} \\
    \small{Dipl\^ome d'Ing\'enieur de Conception (Master 2) : \textbf{Math\'ematiques Appliqu\'ees}, G\'enie Logiciel, Algorithmique.}
\end{itemize}

\section{Comp\'etences Techniques}
\begin{itemize}[leftmargin=*, nosep]
    \item \small{\textbf{MLOps \& Data :} Python, PyTorch, TensorFlow, Scikit-learn, MLflow, pipeline CI/CD, Docker, APIs REST, SQL}
    \item \small{\textbf{ML / Deep Learning :} BERT, Transformers, NLP, classification, extraction d'information, scoring de confiance}
    \item \small{\textbf{Backend \& Infrastructure :} Laravel, Redis, MySQL, Git, Linux, d\'eploiement Cloud}
\end{itemize}

\section{Exp\'eriences \& Projets}
\begin{itemize}[leftmargin=*, label={-}]

\item \textbf{Stage de Recherche -- INRIA Saclay (\'Equipe TRIBE, IP Paris)} \hfill \textit{Avril -- Ao\^ut 2026} \\
\textit{Plateau de Saclay, France} \hfill \textit{Python, BERT, Transformers, NLP, RGPD}
\vspace{-2pt}
    \begin{itemize}[leftmargin=0.4cm, nosep, label=\tiny$\bullet$]
        \item \small{Conception d'un pipeline NLP complet d'analyse de vie priv\'ee sur corpus de 50k+ documents, avec extraction d'entit\'es et scoring de confiance RGPD.}
        \item \small{Fine-tuning de mod\`eles BERT/Transformers pour la classification de clauses contractuelles, atteignant +12\% de F1-score vs baseline.}
        \item \small{Int\'egration du pipeline dans une infrastructure Python Docker pour industrialisation.}
    \end{itemize}
\vspace{3pt}

\item \textbf{Assistant de Recherche -- MILA (Montr\'eal Institute of Learning Algorithms)} \hfill \textit{\'Et\'e 2025} \\
\textit{Collaboration \`a distance} \hfill \textit{PyTorch, TensorFlow, Math\'ematiques}
\vspace{-2pt}
    \begin{itemize}[leftmargin=0.4cm, nosep, label=\tiny$\bullet$]
        \item \small{\'Etude du ph\'enom\`ene Grokking : analyse de la g\'en\'eralisation retard\'ee sur t\^aches arithm\'etiques modulaires avec PyTorch.}
        \item \small{Production de notebooks reproductibles et formalisation des r\'esultats pour publication.}
    \end{itemize}
\vspace{3pt}

\item \textbf{Projet Urban Waste Management -- IA \& IoT} \hfill \textit{2024} \\
\textit{Comp\'etition nationale CITS (1\textsuperscript{er} prix / 75 \'equipes)} \hfill \textit{ML, IoT, Python, classification}
\vspace{-2pt}
    \begin{itemize}[leftmargin=0.4cm, nosep, label=\tiny$\bullet$]
        \item \small{D\'eveloppement d'un syst\`eme de classification automatique des d\'echets via capteurs IoT et mod\`eles ML embarqu\'es.}
        \item \small{Mise en place d'un pipeline de donn\'ees temps r\'eel int\'egrant collecte, pr\'etraitement et inf\'erence.}
    \end{itemize}
\vspace{3pt}

\item \textbf{Projet UROP 2024 -- Analyse ECG par Deep Learning} \hfill \textit{2024} \\
\textit{Collaboration Google DeepMind} \hfill \textit{TensorFlow, Computer Vision, imagerie m\'edicale}
\vspace{-2pt}
    \begin{itemize}[leftmargin=0.4cm, nosep, label=\tiny$\bullet$]
        \item \small{Entra\^inement de CNN pour la d\'etection d'arythmies cardiaques sur signaux ECG, avec augmentation de donn\'ees m\'edicales.}
        \item \small{Optimisation du pipeline d'entra\^inement r\'eduisant le temps de 40\% par vectorisation des pr\'etraitements.}
    \end{itemize}

\end{itemize}

\section{Prix \& Leadership}
\begin{itemize}[leftmargin=*, nosep, label=\tiny$\bullet$]
    \item \small{\textbf{1\textsuperscript{er} prix IndabaX 2025} (IA Sant\'e) ; \textbf{1\textsuperscript{er} prix Mountain Hub 2024} ; \textbf{1\textsuperscript{er} prix CITS National} (75 \'equipes)}
    \item \small{\textbf{Head of Projects Cell \& Lead AI Cell ENSEA} : structuration de 3 p\^oles techniques, +50\% de membres actifs, collaborations Google DeepMind}
    \item \small{\textbf{Open Source :} Laravel API Generator (+30 \'etoiles GitHub, packagist) -- g\'en\'eration de code API via LLM+UML}
\end{itemize}

\end{document}
